% Vertikal tidslinje: dei serielle «fundamenta» (for kapittelet «Dei mislukka
% fundamenta»). Reint TikZ, ingen ekstra bibliotek. Input-bart fragment.
\begin{figure}[t]
\centering
\begin{tikzpicture}[
  yscale=0.92, xscale=1,
  every node/.style={font=\footnotesize},
  aar/.style={font=\small\bfseries},
]
% Loddrett akse
\draw[line width=0.8pt] (0,0.4) -- (0,-16.0);
% Hjelpefunksjon: punkt + år + tekst
\foreach \y/\year/\who/\what in {%
  -0.6/1919/Bauhaus/{Vorkurs: ei grunnlære om form},
  -2.3/1953/Ulm/{vitskap og metode},
  -4.0/1962/Designmetodar/{ein eksplisitt rasjonell prosess},
  -5.7/1969/Simon/{ein vitskap om det kunstige},
  -7.4/1982/Cross/{«designerly ways of knowing»},
  -9.1/1992/Buchanan/{design som det fjerde kunstfaget},
  -10.8/2006/Krippendorff/{«eit nytt fundament for design»},
  -12.5/2012/{Nelson \& Stolterman}/{design som tredje kunnskapskultur},
  -14.2/2017/Redström/{design lagar si eiga teori}%
}{
  \fill (0,\y) circle (2.2pt);
  \node[aar, anchor=east] at (-0.35,\y) {\year};
  \node[anchor=west, text width=78mm] at (0.35,\y) {\textsc{\who} — \textit{\what}};
}
\end{tikzpicture}
\caption{\textsc{Dei serielle fundamenta.} Ni annonserte grunnlag på under hundre år. Eit fag som fann grunnen sin, grunnlegg seg ikkje på nytt kvar generasjon; overfloden er sjølv beviset på fråvêret.}
\end{figure}
